\documentclass[a4paper]{report}

% --- FONT & NGÔN NGỮ ---
\usepackage{fontspec}
\setmainfont{Times New Roman}
\usepackage[vietnamese]{babel}
\usepackage{microtype} % Tự động tinh chỉnh khoảng cách để giảm lỗi tràn lề

% --- ÉP FONT 13pt TOÀN CỤC ---
\usepackage{fontsize}
\changefontsize[15.6]{13} % Cỡ chữ 13pt, base leading 15.6pt

% --- CANH LỀ CHUẨN A4 KHÓA LUẬN TỐT NGHIỆP TDTU ---
% Lề trái 3.5cm (đóng gáy sách), lề phải 2.0cm, lề trên 3.0cm, lề dưới 3.0cm
\usepackage[top=3.0cm, bottom=3.0cm, left=3.5cm, right=2.0cm, headheight=16pt, headsep=1.0cm, footskip=1.2cm]{geometry}

% --- DÃN DÒNG & THỤT ĐẦU DÒNG ĐOẠN VĂN ---
\usepackage{setspace}
\onehalfspacing % Dãn dòng 1.5 lines chuẩn
\usepackage{indentfirst}
\setlength{\parindent}{1.27cm} % Thụt đầu dòng 1.27cm (0.5 inch)
\setlength{\parskip}{0pt}      % Không dãn cách đoạn thừa thãi
\emergencystretch 3em          % Chống tuyệt đối lỗi tràn lề phải (overfull hbox)

% --- HEADER & FOOTER: ĐÁNH SỐ TRANG GIỮA TRÊN ---
\usepackage{fancyhdr}
\pagestyle{fancy}
\fancyhf{} 
\fancyhead[C]{\fontsize{10pt}{12pt}\selectfont\thepage} % Số trang căn giữa phía trên
\renewcommand{\headrulewidth}{0pt}
\renewcommand{\footrulewidth}{0pt}

% Đảm bảo trang đầu chương / trang plain vẫn có số trang căn giữa phía trên
\fancypagestyle{plain}{
    \fancyhf{} 
    \fancyhead[C]{\fontsize{10pt}{12pt}\selectfont\thepage} 
    \renewcommand{\headrulewidth}{0pt}
}

% --- TOÁN HỌC & HÌNH ẢNH ---
\usepackage{amsmath, amsfonts, amssymb}
\usepackage{graphicx}
\usepackage{tikz}
\usetikzlibrary{calc, positioning, decorations.pathreplacing}
\usepackage{float}
\usepackage{longtable}
\usepackage{booktabs}
\usepackage{adjustbox} % bảng sinh tự động: max width=\textwidth (thu nhỏ bảng rộng, không phóng to bảng hẹp)
\usepackage{array}
\usepackage{multirow}
\usepackage{tabularx}
\usepackage{xltabular}
\newcolumntype{Y}{>{\centering\arraybackslash}X}
\usepackage{subcaption}
% --- MÃ GIẢ (THUẬT TOÁN) ---
% Đánh số theo chương (Thuật toán 3.1, ...); từ khóa tiếng Việt; noend: bỏ dòng "hết ..." (thụt lề thể hiện khối lệnh).
\usepackage[chapter]{algorithm}
\usepackage[noend]{algpseudocode}
\floatname{algorithm}{Thuật toán}
\algrenewcommand\algorithmicrequire{\textbf{Đầu vào:}}
\algrenewcommand\algorithmicensure{\textbf{Đầu ra:}}
\algrenewcommand\algorithmicfor{\textbf{với}}
\algrenewcommand\algorithmicforall{\textbf{với mọi}}
\algrenewcommand\algorithmicdo{\textbf{làm}}
\algrenewcommand\algorithmicwhile{\textbf{trong khi}}
\algrenewcommand\algorithmicrepeat{\textbf{lặp}}
\algrenewcommand\algorithmicuntil{\textbf{đến khi}}
\algrenewcommand\algorithmicif{\textbf{nếu}}
\algrenewcommand\algorithmicthen{\textbf{thì}}
\algrenewcommand\algorithmicelse{\textbf{ngược lại}}
\algrenewcommand\algorithmicreturn{\textbf{trả về}}
\algrenewcommand\algorithmicfunction{\textbf{hàm}}
\algrenewcommand\algorithmicprocedure{\textbf{thủ tục}}
\algrenewcommand{\algorithmiccomment}[1]{\hfill\(\triangleright\) \textit{#1}}
\graphicspath{{media/}}
\usepackage{caption}
\captionsetup{
    font={small},
    labelfont={bf},
    labelsep=colon,
    skip=6pt
}

\usepackage[hidelinks]{hyperref} % Liên kết mục lục sạch sẽ, không viền đỏ

% --- TÙY CHỈNH MỤC LỤC & DANH MỤC (TOC, LOF, LOT) ---
\usepackage{tocloft} 
\renewcommand{\cftdotsep}{0.5} % Dấu chấm mục lục san sát chuẩn Word

% Thêm chữ "CHƯƠNG " vào Mục lục
\renewcommand{\cftchappresnum}{CHƯƠNG } 
\renewcommand{\cftchapaftersnum}{. }    
\setlength{\cftchapnumwidth}{6.0em}     

% 1. Cấp Chương (Chapter): In đậm toàn hàng
\renewcommand{\cftchapfont}{\bfseries}
\renewcommand{\cftchapleader}{\bfseries\cftdotfill{\cftdotsep}}
\renewcommand{\cftchappagefont}{\bfseries}

% 2. Cấp Mục (Section - 1.1): Chữ thường
\renewcommand{\cftsecfont}{\normalfont}
\renewcommand{\cftsecleader}{\normalfont\cftdotfill{\cftdotsep}}
\renewcommand{\cftsecpagefont}{\normalfont}
\setlength{\cftsecnumwidth}{2.2em}

% 3. Cấp Tiểu mục (Subsection - 1.1.1): In nghiêng
\renewcommand{\cftsubsecfont}{\itshape}
\renewcommand{\cftsubsecleader}{\itshape\cftdotfill{\cftdotsep}}
\renewcommand{\cftsubsecpagefont}{\itshape}
\setlength{\cftsubsecnumwidth}{2.8em}

% Khoảng cách giữa các dòng trong Mục lục
\setlength{\cftbeforechapskip}{8pt} 
\setlength{\cftbeforesecskip}{5pt}
\setlength{\cftbeforesubsecskip}{4pt}

% Danh mục hình vẽ
\renewcommand{\cftfigpresnum}{Hình } 
\renewcommand{\cftfigaftersnum}{: }  
\setlength{\cftfignumwidth}{4.5em}   

% Danh mục bảng biểu
\renewcommand{\cfttabpresnum}{Bảng }
\renewcommand{\cfttabaftersnum}{: }
\setlength{\cfttabnumwidth}{4.5em}

% Danh mục thuật toán: dòng "Thuật toán 3.1: ..." như danh mục hình/bảng (gói algorithm ghi vào tệp .loa)
\makeatletter
\def\l@algorithm#1#2{{\let\algoldnumberline\numberline
    \renewcommand*{\numberline}[1]{\algoldnumberline{Thuật toán ##1: }}%
    \@dottedtocline{1}{0em}{7.2em}{#1}{#2}}}
\makeatother

% Xóa khoảng trắng thừa giữa các chương trong danh mục hình/bảng
\usepackage{etoolbox}
\makeatletter
\AtBeginDocument{
  \patchcmd{\@chapter}{\addtocontents{lof}{\protect\addvspace{10\p@}}}{}{}{}
  \patchcmd{\@chapter}{\addtocontents{lot}{\protect\addvspace{10\p@}}}{}{}{}
}
\makeatother

% --- ĐỊNH DẠNG TIÊU ĐỀ CÁC CẤP (HEADING SPACING CHUẨN ĐẸP) ---
\usepackage{titlesec}
\setcounter{secnumdepth}{3} % Đánh số tới cấp 1.1.1.1

% Heading 1: Chương (\chapter) - 16pt Bold
\titleformat{\chapter}[hang]
  {\fontsize{16pt}{22pt}\selectfont\bfseries} 
  {CHƯƠNG \thechapter.} 
  {0.5cm} 
  {}
\titlespacing*{\chapter}{0pt}{0pt}{16pt} % Cách dưới 16pt

% Heading 2: Mục lớn (\section) - 14pt Bold, sát lề trái 0cm
\titleformat{\section}[hang]
  {\fontsize{14pt}{19pt}\selectfont\bfseries} 
  {\thesection}
  {0.4cm} 
  {}
\titlespacing*{\section}{0pt}{14pt}{6pt} % Sát lề trái

% Heading 3: Tiểu mục (\subsection - vd 3.1.1) - 13pt Bold Italic, thụt vào 0.635cm (nửa tab)
\titleformat{\subsection}[hang]
  {\fontsize{13pt}{18pt}\selectfont\bfseries\itshape} 
  {\thesubsection}
  {0.35cm} 
  {}
\titlespacing*{\subsection}{0.635cm}{10pt}{4pt} % Thụt vào bậc thang 0.635cm

% Heading 4: Tiểu mục con (\subsubsection - vd 3.1.1.1) - 13pt Italic, thụt vào 1.27cm (1 tab)
\titleformat{\subsubsection}[hang]
  {\fontsize{13pt}{18pt}\selectfont\itshape} 
  {\thesubsubsection}
  {0.3cm} 
  {}
\titlespacing*{\subsubsection}{1.27cm}{8pt}{3pt} % Thụt vào bậc thang 1.27cm

% --- KHOẢNG CÁCH CÔNG THỨC TOÁN HỌC ---
\AtBeginDocument{
  \setlength{\abovedisplayskip}{6pt}
  \setlength{\belowdisplayskip}{6pt}
  \setlength{\abovedisplayshortskip}{3pt} 
  \setlength{\belowdisplayshortskip}{3pt}
}

% --- TÀI LIỆU THAM KHẢO ---
\usepackage{csquotes} 
\DeclareQuoteAlias{english}{vietnamese}

\usepackage[backend=biber, style=ieee]{biblatex}
\addbibresource{backmatter/references.bib}
% biblatex không có gói ngôn ngữ tiếng Việt: thiếu 2 dòng này thì danh mục tài liệu in từ khoá thô
% (and, in, pages, andothers). Dùng chuỗi tiếng Anh, riêng "and"/"et al." đổi sang tiếng Việt.
\DeclareLanguageMapping{vietnamese}{english}
\DefineBibliographyStrings{vietnamese}{andothers={và cộng sự}, and={và}}
\usepackage{xurl} 
\let\parencite\cite 



\defbibheading{bibcenter}{
    \clearpage
    \phantomsection
    \addcontentsline{toc}{chapter}{TÀI LIỆU THAM KHẢO}
    \begin{center}
        {\fontsize{16pt}{22pt}\selectfont\bfseries TÀI LIỆU THAM KHẢO\par}
    \end{center}
    \vspace{12pt}
}

\defbibheading{subbib}{
    \noindent #1\par\vspace{6pt}
}

\defbibenvironment{bibliography}
  {\list
     {}
     {\setlength{\leftmargin}{0pt}
      \setlength{\itemindent}{1.27cm}
      \setlength{\itemsep}{0.5\baselineskip}
      \setlength{\parsep}{0pt}}}
  {\endlist}
  {\item}



% --- TÙY CHỈNH DANH SÁCH (ITEMIZE / ENUMERATE) ---
\usepackage{enumitem}
\setlist{
    topsep=3pt,
    partopsep=0pt,
    itemsep=2pt,
    parsep=0pt,
    labelindent=1.27cm,
    leftmargin=*,
    align=left,
    labelsep=0.4cm
}
\renewcommand{\labelitemi}{--}
\renewcommand{\labelitemii}{$-$}
% --- SỐ LIỆU THỰC NGHIỆM SINH TỰ ĐỘNG ---
% Macro hiện "[chưa có]" nếu số liệu chưa được sinh; bảng/hình chưa có thì hiện khung "Chưa có số liệu".
\providecommand{\resultmacro}[3]{\ifcsname #1-#2-#3\endcsname\csname #1-#2-#3\endcsname\else\textbf{[chưa có]}\fi}
% Giai đoạn phát triển ban đầu (chỉ dùng mã ID; mục Lịch sử phát triển của Chương 4): content/tables/results_macros.tex,
% sinh bởi script đã gỡ khỏi mã nguồn (lịch sử git, commit 9799f1a) -- không còn thay đổi. Mô hình: Pre, Scr, GMF, MLP,
% BPR, MostPop, Random.
\providecommand{\Res}[2]{\resultmacro{res}{#1}{#2}}
\providecommand{\Sig}[2]{\resultmacro{sig}{#1}{#2}}
\InputIfFileExists{content/tables/results_macros.tex}{}{}
% --- SỐ LIỆU H&M (kết quả chính) ---
% Do neumf_project/scripts/24_report_v2.py sinh vào content/tables/v2/ và media/figures/v2/. Không gõ tay số liệu kết quả.
% Mô hình: NeuMFF, LateF, GMFF, MLPF, NeuMF, GMF, MLP, BPR, IKNN, UKNN, Cont, RecPop, Pop, Rand.
% \VRes/\VSd{mô hình}{NDCG|Recall|HR|Precision + K, K = 5, 10, 15, 20}: tập kiểm thử của mẫu kiểm định;
% \VRank{mô hình}{NDCG10}; \VVal: tập xác thực của mẫu phát triển; \VSig{1..10|all}{rel|pholm|lo|hi|verdict|phrase|
% summary}; \VAbl, \VBey, \VData (khóa A = mẫu phát triển, B/Btest = mẫu kiểm định), \VCfg.
\providecommand{\VRes}[2]{\resultmacro{vres}{#1}{#2}}
\providecommand{\VSd}[2]{\resultmacro{vsd}{#1}{#2}}
\providecommand{\VVal}[2]{\resultmacro{vval}{#1}{#2}}
\providecommand{\VRank}[2]{\resultmacro{vrank}{#1}{#2}}
\providecommand{\VSig}[2]{\resultmacro{vsig}{#1}{#2}}
\providecommand{\VAbl}[2]{\resultmacro{vabl}{#1}{#2}}
\providecommand{\VBey}[2]{\resultmacro{vbey}{#1}{#2}}
\providecommand{\VData}[2]{\resultmacro{vdata}{#1}{#2}}
\providecommand{\VCfg}[2]{\resultmacro{vcfg}{#1}{#2}}
\InputIfFileExists{content/tables/v2/results_v2_macros.tex}{}{}
% --- SỐ LIỆU AMAZON BEAUTY (kiểm chứng, audit/PREREG_amazon.md) ---
% Cùng script với NEUMF_CONFIG=configs/amazon.yaml (python run.py amazon-report) sinh vào content/tables/amazon/ và
% media/figures/amazon/. \ARes, \ASd, \ARank, \ASig, \AAbl, \ABey như họ \V...; \AData{all|val|test|run}{...}.
\providecommand{\ARes}[2]{\resultmacro{ares}{#1}{#2}}
\providecommand{\ASd}[2]{\resultmacro{asd}{#1}{#2}}
\providecommand{\ARank}[2]{\resultmacro{arank}{#1}{#2}}
\providecommand{\ASig}[2]{\resultmacro{asig}{#1}{#2}}
\providecommand{\AAbl}[2]{\resultmacro{aabl}{#1}{#2}}
\providecommand{\ABey}[2]{\resultmacro{abey}{#1}{#2}}
\providecommand{\AData}[2]{\resultmacro{adata}{#1}{#2}}
\InputIfFileExists{content/tables/amazon/results_amazon_macros.tex}{}{}
\newcommand{\chuaco}[1]{\begin{center}\fbox{\parbox{0.92\textwidth}{\raggedright\small\textit{Chưa có số liệu: #1}}}\end{center}}
% \bangketquaV{tên file}{lệnh sinh}: chèn bảng sinh tự động trong content/tables/v2/, hoặc khung nhắc nếu chưa có.
\newcommand{\bangketquaV}[2]{\IfFileExists{content/tables/v2/#1}{\input{content/tables/v2/#1}}{\chuaco{chạy #2 rồi \texttt{scripts/24\_report\_v2.py}.}}}
% \bangrefV{tên file}{nhãn}: ``Bảng~X'' nếu bảng đã được sinh, nếu chưa thì chữ thay thế (tránh tham chiếu ``??'').
\newcommand{\bangrefV}[2]{\IfFileExists{content/tables/v2/#1}{Bảng~\ref{#2}}{bảng kết quả (chưa có số liệu)}}
% \bangketquaA, \bangrefA: như trên cho bảng Amazon Beauty trong content/tables/amazon/.
\newcommand{\bangketquaA}[1]{\IfFileExists{content/tables/amazon/#1}{\input{content/tables/amazon/#1}}{\chuaco{chạy \texttt{python run.py amazon-report}.}}}
\newcommand{\bangrefA}[2]{\IfFileExists{content/tables/amazon/#1}{Bảng~\ref{#2}}{bảng kết quả Amazon (chưa có số liệu)}}
% \hinhref{đường dẫn hình}{nhãn}: ``Hình~X'' nếu hình đã được sinh, nếu chưa thì chữ thay thế.
\newcommand{\hinhref}[2]{\IfFileExists{#1}{Hình~\ref{#2}}{hình kết quả (chưa có số liệu)}}
% \hinhketqua{đường dẫn}{độ rộng}{chú thích}{nhãn}{lệnh sinh ra hình}
\newcommand{\hinhketqua}[5]{\IfFileExists{#1}{\begin{figure}[htbp]\centering\includegraphics[width=#2]{#1}\caption{#3}\label{#4}\end{figure}}{\chuaco{hình ``#3'' --- #5.}}}
